\documentclass[12pt]{article}

\usepackage[margin=1in]{geometry}
\usepackage{amsmath,amssymb}
\usepackage{setspace}
\usepackage{booktabs}
\usepackage{natbib}
\usepackage{hyperref}
\usepackage{graphicx}

\hypersetup{
    colorlinks=true,
    linkcolor=blue,
    citecolor=blue,
    urlcolor=blue
}

\onehalfspacing

\title{From State Houses to City Halls:\\
The Devolution of Government Finance\\
in Nineteenth-Century America}

\author{Notes compiled by Thomas J.\ Sargent and Claude\\[6pt]
\small Drawing on the work of John Joseph Wallis and collaborators}

\date{March 2026}

\begin{document}
\maketitle

\begin{abstract}
During the nineteenth century, the locus of subnational government finance in the United States shifted dramatically---from state governments, which dominated fiscal activity in the early Republic, to municipalities and counties, which by 1900 collected and spent far more than their parent states.   This note surveys the evidence for this transformation, drawing heavily on the research of John Joseph Wallis and his collaborators.  We emphasize the role of the state debt crises of the early 1840s, the constitutional reforms that followed, and the consequent devolution of taxing and spending authority to local governments.
A key quantitative marker: by the time of the first Census of Governments in 1902, local government revenues exceeded state revenues by a ratio of roughly four to one.
\end{abstract}

\tableofcontents

\newpage

\section{Introduction}

A striking feature of American fiscal history is that  states and municipalities reversed their relative fiscal  positions  during the nineteenth century.  In the early decades of the Republic, state governments were the dominant subnational fiscal actors.  They collected property taxes, earned revenue from state-owned banks and investments, and, most consequentially, borrowed heavily to finance canals, turnpikes, and railroads.  Local governments---counties, towns, cities---played a subordinate role, administering justice, maintaining local roads, and supporting the poor.

By the end of the century, this relationship had been inverted.  Cities like New York, Philadelphia, Chicago, and Boston operated budgets that dwarfed those of their respective state governments.  Municipal governments built and operated waterworks, sewage systems, street lighting, professional police and fire departments, public schools, parks, and libraries.  County governments ran courts, maintained rural roads, and administered property tax collections that fed both county and municipal treasuries.

John Joseph Wallis of the University of Maryland has documented this transformation more thoroughly than any other scholar.  His contributions include the chapter on ``Government Finances and Employment'' in the \textit{Historical Statistics of the United States} (Cambridge University Press, 2006), his \textit{Journal of Economic Perspectives} article ``American Government Finance in the Long Run: 1790 to 1990'' \citep{wallis2000jep}, his chapter ``A History of the Property Tax in America'' \citep{wallis2001property}, his article with Legler and Sylla on ``U.S.\ City Finances and the Growth of Government, 1850--1902'' \citep{legler_sylla_wallis1988}, and his work with Wallace Oates on ``Decentralization in the Public Sector'' \citep{wallis_oates1988}.

This note synthesizes the evidence, organizes the key statistics, and discusses the causal mechanisms that drove the shift.


\section{The Fiscal Landscape of the Early Republic, 1790--1840}

\subsection{Federal government: tariffs and land sales}

Through most of the antebellum period, the federal government relied on two revenue sources: customs duties (tariffs on imported goods) and proceeds from sales of public lands.  Together these revenues were more than sufficient to fund the limited activities of the national government---defense, diplomacy, postal service, and servicing the national debt.  The federal government accounted for roughly 55--65\% of total government revenue in the early Republic; this share would decline by century's end \citep{wallis2000jep}.

\subsection{State governments: the dominant subnational fiscal actors}

In the 1820s and 1830s, state governments were the principal subnational fiscal actors.  Their revenues came from several sources:

\begin{itemize}
    \item \textbf{Property taxes.}  Most states imposed a general property tax, often assessed and collected by county officials but authorized and controlled by the state legislature.  The property tax was initially a \textit{state} tax, not a local one \citep{wallis2001property}.

    \item \textbf{Revenue from state-owned enterprises and investments.}  Several states owned or invested in banks, receiving dividend income.  Some states earned tolls from canals and turnpikes.  Others received revenue from land sales.

    \item \textbf{Bank taxes and charter fees.}  States taxed banks and granted corporate charters for fees, creating a significant revenue stream, particularly in states like Pennsylvania and New York \citep{sylla_legler_wallis1987}.

    \item \textbf{Borrowing.}  States borrowed aggressively in the 1830s to finance ``internal improvements''---canals, railroads, and turnpikes.  Total state debt grew from approximately \$13 million in 1825 to nearly \$200 million by 1841 \citep{wallis2005constitutions, grinath1997debt}.
\end{itemize}

\subsection{Local governments: modest and subordinate}

Local governments in this period had limited fiscal responsibilities.  Counties administered justice, maintained local roads (often through labor service requirements rather than tax-financed expenditure), and supported the poor through ``poor houses'' and ``outdoor relief.''  Towns and cities provided rudimentary policing (often through night watches), maintained streets, and regulated local markets.

Local revenues were correspondingly small.  Cities and counties collected property taxes, but the rates were low and the tax base was limited.  Many local functions were financed through fees, fines, and in-kind labor obligations rather than taxation.

Total local government expenditures in 1840 were modest compared to state expenditures.  Wallis's estimates suggest that state expenditures were roughly comparable to, and in many states exceeded, local expenditures in the 1830s \citep{wallis2000jep}.


\section{The Crisis of the 1840s: A Fiscal Watershed}

\subsection{State defaults and repudiation}

The financial crisis that followed the Panic of 1837 precipitated a wave of state defaults that fundamentally altered the trajectory of American public finance.  Between 1841 and 1843, nine states defaulted on their debts: Pennsylvania, Maryland, Indiana, Illinois, Michigan, Arkansas, Mississippi, Louisiana, and Florida.  Four of these---Mississippi, Florida, Arkansas, and Michigan---ultimately repudiated all or part of their debts \citep{wallis2004sovereign, grinath1997debt}.

The defaults were concentrated among states that had borrowed heavily to finance internal improvements using what \citet{wallis2005constitutions} terms ``taxless finance''---borrowing against anticipated toll revenues and rising land values rather than raising taxes to cover costs.

\subsection{Constitutional reforms}

The defaults triggered a wave of state constitutional reforms during the 1840s and early 1850s.  These reforms shared several common features:

\begin{enumerate}
    \item \textbf{Debt limitations}: States imposed constitutional ceilings on state debt, typically requiring popular referenda or supermajority legislative votes for any borrowing above a specified threshold.

    \item \textbf{Balanced budget requirements}: Many states required that taxes be levied simultaneously with appropriations, eliminating the option of ``taxless finance.''

    \item \textbf{Prohibitions on state investment in private enterprises}: States were forbidden from purchasing stock in, or lending credit to, private corporations.

    \item \textbf{Prohibitions on assumption of local debts}: States were prohibited from assuming the debts of local governments or private corporations.

    \item \textbf{General incorporation laws}: States replaced the system of special legislative charters with general incorporation statutes, reducing opportunities for legislative corruption \citep{wallis2005constitutions}.
\end{enumerate}

These constitutional reforms had far-reaching consequences.  By constraining state borrowing and investment, they dramatically reduced the fiscal role of state governments.  But the demand for infrastructure investment did not disappear---it devolved to local governments.


\section{The Devolution to Local Government, 1850--1900}

\subsection{The mechanism: what states abandoned, cities assumed}

The constitutional reforms of the 1840s created what \citet{walliswein2008} describe as a deliberate institutional architecture rather than an unintended consequence.  By tying state legislators' hands, the reforms forced infrastructure investment decisions down to the local level, where residents who would benefit from improvements could vote on (and tax themselves for) specific projects.

Cities and counties stepped into the vacuum left by states.  Municipal governments undertook the construction and operation of:

\begin{itemize}
    \item Waterworks and sewage systems
    \item Paved streets and bridges
    \item Professional police and fire departments
    \item Public school systems
    \item Parks, libraries, and public health infrastructure
    \item Gas and later electric lighting
    \item Street railways and transit systems (often through franchise agreements)
\end{itemize}

Counties retained their traditional functions (courts, roads, poor relief) and in many states assumed responsibility for rural school finance.

\subsection{The property tax becomes a local tax}

Perhaps the single most important institutional change was the devolution of the property tax.  Before 1840, the general property tax was primarily a state-level revenue instrument.  County officials assessed and collected the tax, but the legislature determined its rate and claimed all or most of its proceeds.

After 1840, states increasingly ceded the property tax to local governments.  Several forces drove this transformation:

\begin{enumerate}
    \item \textbf{State debt constraints}: With borrowing limited, states needed less revenue and were content to leave the property tax to localities.

    \item \textbf{Local spending needs}: Urbanization created demands for expensive infrastructure that required substantial tax revenue at the local level.

    \item \textbf{Administrative logic}: The property tax was inherently local in its administration---assessors needed local knowledge to value real estate---and it made sense for the government that administered the tax to also receive the revenue.

    \item \textbf{Alternative state revenue sources}: States increasingly turned to corporate taxes, inheritance taxes, and fees to replace property tax revenue \citep{wallis2001property}.
\end{enumerate}

By the late nineteenth century, the property tax had become overwhelmingly a local tax.  In 1902, local governments collected approximately \$624 million in property taxes, while state governments collected roughly \$82 million.  The property tax accounted for about 73\% of all local government revenue but only about 40\% of state government revenue.

\subsection{Urbanization as a fiscal driver}

Urbanization both drove and was driven by the growth of local government.  The urban population of the United States rose from approximately 1.8 million (10.8\% of the total) in 1840 to 30.4 million (39.7\%) in 1900.  Cities that barely existed in 1840---Chicago, San Francisco, Minneapolis, Denver---had populations in the hundreds of thousands by 1900.

Urban density created both the need for collective provision of water, sewage, fire protection, and policing, and a concentrated tax base of improved real property that could finance these services.  As \citet{legler_sylla_wallis1988} document, city expenditures per capita grew substantially between 1850 and 1902, driven by capital investment in urban infrastructure.


\section{The Evidence: Revenue and Expenditure Statistics}

\subsection{Orders of magnitude}

The quantitative evidence for the state-to-local shift is striking.  Table~\ref{tab:rev_comparison} summarizes the approximate distribution of government revenues across the three levels of government at selected dates during the nineteenth century, drawing on figures assembled by Wallis for the \textit{Historical Statistics of the United States} and the \textit{Journal of Economic Perspectives}.

\begin{table}[htbp]
\centering
\caption{Distribution of Government Revenue by Level of Government (\%)}
\label{tab:rev_comparison}
\begin{tabular}{@{}lccc@{}}
\toprule
Year & Federal & State & Local \\
\midrule
1840 & 62 & 21 & 17 \\
1860 & 65 & 12 & 23 \\
1870 & 68 & 10 & 22 \\
1880 & 55 & 11 & 34 \\
1890 & 49 & 11 & 40 \\
1902 & 38 & 12 & 50 \\
\bottomrule
\end{tabular}
\smallskip

\small\textit{Sources}: Approximate shares based on figures assembled by \citet{wallis2000jep} and the chapter on government finances in the \textit{Historical Statistics of the United States} \citep{wallis2006hsus}.  Figures for years before 1902 involve some estimation and interpolation.  Civil War years distort the federal share in 1870.
\end{table}

\noindent Several features of Table~\ref{tab:rev_comparison} stand out:
\begin{itemize}
    \item The local share more than doubled between 1840 and 1902, rising from roughly 17\% to 50\%.
    \item The state share actually declined from about 21\% to 12\% between 1840 and 1902.
    \item The federal share, though dominant at mid-century (boosted by tariff revenue), declined as local government grew.
\end{itemize}

\subsection{The 1902 Census of Governments}

The first comprehensive Census of Governments in 1902 provides the most reliable snapshot of the fiscal landscape at the end of the transformation.  The key figures show the overwhelming dominance of local government:

\begin{table}[htbp]
\centering
\caption{Government Revenues, 1902 (millions of dollars)}
\label{tab:1902}
\begin{tabular}{@{}lrc@{}}
\toprule
Level of government & Revenue & Share \\
\midrule
Federal & 562 & 38.2\% \\
State & 183 & 12.4\% \\
Local & 726 & 49.4\% \\
\midrule
Total & 1{,}471 & 100.0\% \\
\bottomrule
\end{tabular}
\smallskip

\small\textit{Source}: Bureau of the Census, \textit{Wealth, Debt, and Taxation: 1902}; figures as tabulated by \citet{wallis2000jep}.
\end{table}

\noindent Local government revenue in 1902 was approximately four times that of state governments.  Within local government, cities (incorporated municipalities) accounted for the largest share, followed by counties and school districts.

\subsection{The state-to-local ratio over time}

A useful summary statistic is the ratio of local to state revenue (or expenditure).  This ratio captures the essence of the transformation:

\begin{table}[htbp]
\centering
\caption{Approximate Ratio of Local to State Revenue}
\label{tab:ratio}
\begin{tabular}{@{}lc@{}}
\toprule
Year & Local/State ratio \\
\midrule
1840 & $\approx 0.8$ \\
1860 & $\approx 1.9$ \\
1880 & $\approx 3.1$ \\
1902 & $\approx 4.0$ \\
\bottomrule
\end{tabular}
\smallskip

\small\textit{Source}: Computed from estimates in \citet{wallis2000jep} and \citet{wallis2006hsus}.
\end{table}

\noindent The ratio rose from less than one (states collected more) around 1840 to approximately four (local governments collected four times as much) by 1902.

\subsection{Property tax collections by level of government}

The devolution of the property tax is particularly well documented:

\begin{table}[htbp]
\centering
\caption{Property Tax Revenue by Level of Government, 1902}
\label{tab:proptax}
\begin{tabular}{@{}lrc@{}}
\toprule
Level & Property tax revenue & Share \\
\midrule
State & \$82 million & 11.6\% \\
Local & \$624 million & 88.4\% \\
\midrule
Total & \$706 million & 100.0\% \\
\bottomrule
\end{tabular}
\smallskip

\small\textit{Source}: Bureau of the Census, \textit{Wealth, Debt, and Taxation: 1902}; \citet{wallis2001property}.
\end{table}

\noindent By 1902, local governments collected  88\% of all property tax revenue.  Several states had by then abandoned the state property tax entirely, relying instead on corporate taxes, inheritance taxes, license fees, and other revenue sources.  The complete separation of the property tax from the state level would continue into the twentieth century: by the 1930s, most states had abandoned the general property tax altogether \citep{wallis2001property}.


\section{City Finance: The Engine of Growth}

\subsection{The Legler--Sylla--Wallis findings}

\citet{legler_sylla_wallis1988} provide the most detailed study of city finance in the late nineteenth century.  Using data from the Census on cities with populations over 4,000, they document several important patterns:

\begin{enumerate}
    \item \textbf{Rapid absolute growth}: Total city revenues grew from approximately \$48 million in 1860 to \$640 million in 1902, a more than thirteen-fold increase.  After adjusting for price changes (which were modest over this period due to gold standard deflation), real growth was still roughly eight-fold.

    \item \textbf{Per capita growth}: Real per capita city expenditures roughly tripled between 1860 and 1902, even as the urban population grew enormously.

    \item \textbf{Capital-intensive spending}: A large share of city expenditure went to capital investment in infrastructure---waterworks, sewers, streets, bridges, and public buildings.  Cities borrowed heavily to finance these investments, with municipal debt growing from about \$100 million in 1860 to over \$1.5 billion by 1902.

    \item \textbf{The property tax as mainstay}: The general property tax on real estate provided the foundation of city revenue, supplemented by license fees, special assessments (levied on property owners who directly benefited from improvements like street paving), and payments from franchised utilities.
\end{enumerate}

\subsection{What cities spent money on}

By the late nineteenth century, the expenditure pattern of a typical large American city looked roughly as follows:

\begin{itemize}
    \item \textbf{Streets, highways, and bridges}: 15--20\% of total expenditure
    \item \textbf{Police}: 8--12\%
    \item \textbf{Fire}: 5--8\%
    \item \textbf{Public education}: 15--25\% (growing over time)
    \item \textbf{Water supply and sewerage}: 10--15\%
    \item \textbf{Debt service}: 15--25\% (reflecting heavy past borrowing for infrastructure)
    \item \textbf{Public health and charities}: 5--10\%
    \item \textbf{Parks and recreation}: 2--5\%
    \item \textbf{General government}: 5--8\%
\end{itemize}

This expenditure pattern reveals the essentially local character of government in late-nineteenth-century America.  The services that most directly affected citizens' daily lives---clean water, safe streets, fire protection, schooling---were delivered and financed at the municipal level.


\section{Why Did the Shift Occur? Interpretive Frameworks}

\subsection{The constitutional-reform channel}

The most direct explanation, emphasized by \citet{wallis2005constitutions}, is institutional.  The state constitutional reforms of the 1840s deliberately constrained state fiscal activity.  Debt ceilings, balanced budget requirements, and prohibitions on state investment in private enterprises all reduced the scope for state-level fiscal action.  These constraints did not kill the demand for infrastructure; they redirected it to the local level.

In this view, the shift was not primarily driven by urbanization per se, though urbanization magnified its effects.  It was driven by a deliberate political choice to restructure American federalism after the traumatic experience of state defaults.

\subsection{The urbanization channel}

A complementary explanation emphasizes the demand side.  Rapid urbanization after 1840 created enormous demand for collective goods---water, sewerage, fire protection, policing---that only local governments could efficiently provide.  Rural areas required fewer public services and lower per capita expenditure.  As the population shifted from rural to urban settings, total local government spending naturally rose.

\subsection{The public-choice interpretation: accountability and fiscal discipline}

\citet{walliswein2008} offer a public-choice interpretation of the post-1840 institutional structure.  They argue that the devolution of spending authority to the local level was not ``dysfunctional'' but rather \textit{optimal} in a second-best sense.  By tying specific expenditures to specific local tax increases, the post-reform structure forced greater accountability: residents who would benefit from a new waterworks or school had to vote to tax themselves to pay for it.

This contrasts with the pre-1840 regime, in which state legislators could approve expensive canal projects financed by borrowing, with the costs diffused across the entire state and deferred to future taxpayers.  The devolution to local government was thus a mechanism for improving the alignment between fiscal costs and benefits.

\subsection{The property tax as a coordinating device}

In a striking case study, \citet{wallis2003indiana} examines how Indiana used the property tax to coordinate the financing of its ``Mammoth System'' of internal improvements in the 1835--1842 period.  The property tax served as a device for allocating costs across benefiting districts.  When the system collapsed, the lesson learned was that investment decisions should be made and financed locally---a principle that the post-1840 constitutional reforms enshrined.


\section{The Declining Fiscal Role of State Governments}

State governments did not disappear from fiscal life, but their role changed qualitatively as well as quantitatively.

\subsection{What states continued to do}

After the constitutional reforms, state governments focused on:

\begin{itemize}
    \item \textbf{Maintaining courts and the justice system}
    \item \textbf{Operating prisons and asylums}
    \item \textbf{Chartering and regulating corporations, banks, and railroads}
    \item \textbf{Supporting limited education activities} (normal schools, state universities)
    \item \textbf{Minimal regulatory functions} (weights and measures, militia)
\end{itemize}

These functions required modest budgets.  The typical state government of the 1870s--1890s spent relatively little compared to the dynamic and growing cities within its borders.

\subsection{State revenue diversification}

As states ceded the property tax to localities, they sought alternative revenue sources.  The transition unfolded gradually over the second half of the nineteenth century and into the twentieth:

\begin{itemize}
    \item \textbf{Corporate franchise taxes}: Taxes on the privilege of incorporation became increasingly important, especially in states like New Jersey and later Delaware.
    \item \textbf{Inheritance taxes}: Several states adopted inheritance taxes in the late nineteenth century.
    \item \textbf{License fees}: States imposed fees for various licenses and permits.
    \item \textbf{Insurance company taxes}: Taxes on premiums of insurance companies operating within the state.
    \item \textbf{Selective sales taxes}: Taxes on specific commodities like alcohol and tobacco.
\end{itemize}

It was not until the twentieth century---with the adoption of state income taxes (Wisconsin in 1911 being the first modern version) and general sales taxes (West Virginia in 1921)---that states developed the revenue instruments that would allow them to reassert fiscal significance.


\section{Connections to the State Debt Crisis Literature}

This note complements our earlier work on the state defaults of the 1840s.  The key connection is causal: the defaults triggered the constitutional reforms \citep{wallis2005constitutions}, which constrained state fiscal activity, which in turn drove the devolution of spending and taxing authority to local governments.

The chain of causation can be summarized as:
\[
\text{Taxless finance} \rightarrow \text{Defaults (1841--43)} \rightarrow \text{Constitutional reforms (1840s--50s)} \rightarrow \text{State fiscal retrenchment} \rightarrow \text{Local government expansion}
\]

This sequence makes the state defaults of the 1840s not merely an episode in financial history, but a pivotal event in the long-run evolution of American fiscal federalism.


\section{Summary and Implications}

\begin{enumerate}
    \item In the early Republic (1790--1840), state governments were the dominant subnational fiscal actors.  They collected the bulk of property taxes, earned revenue from banks and investments, and borrowed heavily to finance internal improvements.

    \item The state debt crises of the early 1840s triggered constitutional reforms that sharply constrained state borrowing and spending.

    \item After 1840, local governments---especially municipalities---assumed the primary role in infrastructure investment and service provision.  This shift was driven both by institutional changes (constitutional debt limits) and by urbanization (which created demand for local public goods).

    \item By 1902, local government revenues were approximately four times state revenues.  The property tax, once primarily a state revenue instrument, had become overwhelmingly a local tax.

    \item The shift from state to local fiscal dominance was one of the most important structural changes in American public finance during the nineteenth century.  It shaped the decentralized, locally financed character of American government that persists to this day.
\end{enumerate}

As John Wallis has emphasized throughout his work, understanding the fiscal history of the United States requires attending not just to the federal level---where most historical attention has focused---but to the state and local levels, where most of the growth and innovation in government actually occurred during the nineteenth century.

\newpage

\begin{thebibliography}{99}

\bibitem[Grinath et~al.(1997)]{grinath1997debt}
Grinath, Arthur, John Joseph Wallis, and Richard Sylla.
\newblock ``Debt, Default, and Revenue Structure: The American State Debt Crisis in the Early 1840s.''
\newblock NBER Historical Working Paper No.\ 97, 1997.

\bibitem[Legler et~al.(1988)]{legler_sylla_wallis1988}
Legler, John B., Richard Sylla, and John Joseph Wallis.
\newblock ``U.S.\ City Finances and the Growth of Government, 1850--1902.''
\newblock \textit{Journal of Economic History} 48(2): 347--356, 1988.

\bibitem[Sylla et~al.(1987)]{sylla_legler_wallis1987}
Sylla, Richard, John B.\ Legler, and John Joseph Wallis.
\newblock ``Banks and State Public Finance in the New Republic: The United States, 1790--1860.''
\newblock \textit{Journal of Economic History} 47(2): 391--403, 1987.

\bibitem[Wallis(2000)]{wallis2000jep}
Wallis, John Joseph.
\newblock ``American Government Finance in the Long Run: 1790 to 1990.''
\newblock \textit{Journal of Economic Perspectives} 14(1): 61--82, 2000.

\bibitem[Wallis(2001)]{wallis2001property}
Wallis, John Joseph.
\newblock ``A History of the Property Tax in America.''
\newblock In Wallace E.\ Oates, ed., \textit{Property Taxation and Local Government Finance}, pp.\ 123--147. Cambridge: Lincoln Institute of Land Policy, 2001.

\bibitem[Wallis(2003)]{wallis2003indiana}
Wallis, John Joseph.
\newblock ``The Property Tax as a Coordinating Device: Financing Indiana's Mammoth System of Internal Improvements, 1835 to 1842.''
\newblock \textit{Explorations in Economic History} 40(3): 223--250, 2003.

\bibitem[Wallis(2005)]{wallis2005constitutions}
Wallis, John Joseph.
\newblock ``Constitutions, Corporations, and Corruption: American States and Constitutional Change, 1842--1852.''
\newblock \textit{Journal of Economic History} 65(1): 211--256, 2005.

\bibitem[Wallis(2006)]{wallis2006hsus}
Wallis, John Joseph.
\newblock ``Government Finances and Employment.''
\newblock Chapter Ea in Susan B.\ Carter et~al., eds., \textit{Historical Statistics of the United States: Earliest Times to the Present}, Millennial Edition. New York: Cambridge University Press, 2006.

\bibitem[Wallis(2007)]{wallis2007promotion}
Wallis, John Joseph.
\newblock ``American Government and the Promotion of Economic Development in the National Era, 1790 to 1860.''
\newblock In Price Fishback, ed., \textit{Government and the American Economy: A New History}, pp.\ 148--185. Chicago: University of Chicago Press, 2007.

\bibitem[Wallis and Oates(1988)]{wallis_oates1988}
Wallis, John Joseph, and Wallace E.\ Oates.
\newblock ``Decentralization in the Public Sector: An Empirical Study of State and Local Government.''
\newblock In Harvey S.\ Rosen, ed., \textit{Fiscal Federalism: Quantitative Studies}, pp.\ 5--32. Chicago: University of Chicago Press, 1988.

\bibitem[Wallis and Weingast(2008)]{walliswein2008}
Wallis, John Joseph, and Barry R.\ Weingast.
\newblock ``Dysfunctional or Optimal Institutions: State Debt Limitations, the Structure of State and Local Governments, and the Finance of American Infrastructure.''
\newblock In Elizabeth Garrett, Elizabeth A.\ Graddy, and Howell E.\ Jackson, eds., \textit{Fiscal Challenges: An Interdisciplinary Approach to Budget Policy}, pp.\ 331--365. New York: Cambridge University Press, 2008.

\bibitem[Wallis et~al.(2004)]{wallis2004sovereign}
Wallis, John Joseph, Richard Sylla, and Arthur Grinath.
\newblock ``Sovereign Debt and Repudiation: The Emerging-Market Debt Crisis in the U.S.\ States, 1839--1843.''
\newblock NBER Working Paper No.\ 10753, 2004.

\end{thebibliography}

\end{document}
